% Unit 094 terjemahan Bahasa Indonesia dari chapter7.tex baris 1280--1444.
% Sumber: Methods of Algebra, Volume 2, commit
% 9a5803ff77dd3257484cb177f851a73770a59dd3 (CC BY 4.0).
% stable-unit-id: o014.aljabr2.chapter7.beck-monadicity-a
% Cakupan: Bagian 7.6 dari pembukaan hingga konstruksi modul-T bebas;
% dilanjutkan oleh Unit 095 pada baris 1445.

% segment-id: o014.li2.chapter7.beck-monadicity-a.h001
\section{Teorema Monadisitas Beck}\label{sec:Beck}
\hypertarget{o014.aljabr2.chapter7.beck-monadicity-a}{ }
% segment-id: o014.li2.chapter7.beck-monadicity-a.p001
Misalkan $\mathcal{C}$ kategori. Tinjau kategori funktor
$\mathcal{C}^{\mathcal{C}}$ yang objek-objeknya semua endofunktor
$T:\mathcal{C}\to\mathcal{C}$ dan morfisme-morfismenya morfisme
antarendofunktor. Kategori ini secara alami menjadi kategori monoidal: perkalian
$\otimes$ adalah komposisi funktor dan objek satuan $\munit$ adalah funktor
identitas. Karena komposisi funktor asosiatif secara ketat, kategori ini juga
merupakan kategori monoidal ketat.

% segment-id: o014.li2.chapter7.beck-monadicity-a.d001
\begin{definition}[Monad dan komonad]\label{def:monad}
	\index{monad@monad}
	\index{komonad@komonad (comonad)}
	Misalkan $\mathcal{C}$ kategori. Aljabar dalam kategori monoidal
	$\mathcal{C}^{\mathcal{C}}$ (Definisi~\ref{def:algebra-monoidal}) disebut
	\emph{monad} di atas $\mathcal{C}$. Secara dual, monad di atas
	$\mathcal{C}^{\opp}$ disebut \emph{komonad} di atas $\mathcal{C}$.
\end{definition}

% segment-id: o014.li2.chapter7.beck-monadicity-a.p002
Dengan menguraikan Definisi~\ref{def:algebra-monoidal} dalam kategori
endofunktor $\mathcal{C}^{\mathcal{C}}$, diperoleh bahwa monad adalah data
$(T,\mu,\eta)$, dengan $T:\mathcal{C}\to\mathcal{C}$ suatu funktor serta
$\mu:T^2\to T$ dan $\eta:\identity_{\mathcal{C}}\to T$ suatu morfisme,
sedemikian sehingga diagram berikut komutatif:
% segment-id: o014.li2.chapter7.beck-monadicity-a.q001
\begin{equation*}\begin{gathered}
	\begin{tikzcd}
		T \arrow[r, "{\eta T}"] \arrow[rd, "\identity"'] & T^2 \arrow[d, "\mu"] & T \arrow[l, "{T\eta}"'] \arrow[ld, "\identity"] \\
		& T &
	\end{tikzcd} \quad
	\begin{tikzcd}
		T^3 \arrow[r, "\mu T"] \arrow[d, "T\mu"'] & T^2 \arrow[d, "\mu"] \\
		T^2 \arrow[r, "\mu"'] & T.
	\end{tikzcd}
\end{gathered}\end{equation*}

% segment-id: o014.li2.chapter7.beck-monadicity-a.p003
Dibandingkan dengan definisi semula, kendala asosiativitas tidak lagi
diperlukan di sini, sedangkan $\eta\otimes\identity$ (atau
$\identity\otimes\eta$) diterjemahkan menjadi $\eta T$ (atau $T\eta$), dan
demikian seterusnya.

% segment-id: o014.li2.chapter7.beck-monadicity-a.p004
Secara dual, komonad $(L,\delta,\epsilon)$ terdiri atas funktor
$L:\mathcal{C}\to\mathcal{C}$ serta morfisme $\delta:L\to L^2$ dan
$\epsilon:L\to\identity_{\mathcal{C}}$, dengan syarat bahwa diagram berikut
komutatif.
% segment-id: o014.li2.chapter7.beck-monadicity-a.q002
\begin{equation*}\begin{gathered}
	\begin{tikzcd}
		L  & L^2 \arrow[l, "{\epsilon L}"'] \arrow[r, "{L\epsilon}"] & L \\
		& L \arrow[lu, "\identity"] \arrow[ru, "\identity"'] \arrow[u, "\delta"] &
	\end{tikzcd} \quad
	\begin{tikzcd}
		L^3 & L^2 \arrow[l, "\delta L"'] \\
		L^2 \arrow[u, "L\delta"] & L \arrow[u, "\delta"'] \arrow[l, "\delta"]
	\end{tikzcd}
\end{gathered}\end{equation*}

% segment-id: o014.li2.chapter7.beck-monadicity-a.p005
Data $(T,\mu,\eta)$ (atau $(L,\delta,\epsilon)$) lazim disingkat menjadi $T$
(atau $L$). Sumber utama monad dan komonad adalah pasangan adjoin.

% segment-id: o014.li2.chapter7.beck-monadicity-a.eg001
\begin{example}[Pasangan adjoin menentukan monad]\label{eg:adjunction-monad}
	Tinjau sepasang funktor adjoin
	% segment-id: o014.li2.chapter7.beck-monadicity-a.q003
	\[\begin{tikzcd}
		F: \mathcal{C} \arrow[shift left, r] & \mathcal{D}: G \arrow[shift left, l],
	\end{tikzcd}\]
	beserta unit $\eta:\identity_{\mathcal{C}}\to GF$ dan kounit
	$\varepsilon:FG\to\identity_{\mathcal{D}}$ yang bersesuaian. Definisikan
	monad $(T,\mu,\eta)$ di atas $\mathcal{C}$ dan komonad
	$(L,\delta,\varepsilon)$ di atas $\mathcal{D}$ sebagai berikut.
	% segment-id: o014.li2.chapter7.beck-monadicity-a.q004
	\[\begin{array}{|c|c|}\hline
		T := GF & L := FG \\
		\mu := \left[ GFGF \xrightarrow{G\varepsilon F} GF \right] & \delta := \left[ FG \xrightarrow{F\eta G} FGFG \right] \\
		\eta : \identity_{\mathcal{C}} \to GF & \varepsilon: FG \to \identity_{\mathcal{D}} \\ \hline
	\end{array}\]

	Berdasarkan dualitas, cukup memeriksa aksioma monad dalam kasus
	$(T,\mu,\eta)$. Pertama, diagram untuk unit ialah
	% segment-id: o014.li2.chapter7.beck-monadicity-a.q005
	\[\begin{tikzcd}[column sep=large]
		GF \arrow[r, "{\eta GF}"] \arrow[equal, rd] & GFGF \arrow[d, "{G\varepsilon F}"] & GF \arrow[l, "{GF\eta}"'] \arrow[equal, ld] \\
		& GF &
	\end{tikzcd}\]
	yang komutatif berkat identitas segitiga
	$(G\varepsilon)(\eta G)=\identity_G$ dan
	$(\varepsilon F)(F\eta)=\identity_F$. Diagram asosiativitas ialah
	% segment-id: o014.li2.chapter7.beck-monadicity-a.q006
	\[\begin{tikzcd}[column sep=large]
		GFGFGF \arrow[d, "{G\varepsilon FGF}"'] \arrow[r, "{GFG\varepsilon F}"] & GFGF \arrow[d, "G\varepsilon F"] \\
		GFGF \arrow[r, "G\varepsilon F"'] & GF
	\end{tikzcd}\]
	dan diagram ini komutatif: kedua komposisinya sama-sama digambarkan sebagai
	% segment-id: o014.li2.chapter7.beck-monadicity-a.q007
	\[\begin{tikzcd}
		\mathcal{C} \arrow[r, "F"] & \mathcal{D} \arrow[r, "G"] \arrow[rr, bend right=60, ""{name=A}, "{\identity}"'] & \mathcal{C} \arrow[r, "F"] \arrow[Rightarrow, to=A, "\varepsilon"] &
		\mathcal{D} \arrow[r, "G"] \arrow[rr, bend right=60, ""{name=B}, "{\identity}"'] & \mathcal{C} \arrow[r, "F"] \arrow[Rightarrow, to=B, "\varepsilon"] & \mathcal{D} \arrow[r, "G"] & \mathcal{C}
	\end{tikzcd}\]
	dengan hanya urutan komposisi dua daerah melengkung yang berbeda, sedangkan
	hasilnya sama. Ini merupakan kasus khusus hukum pertukaran antara komposisi
	vertikal dan horizontal morfisme \cite[Lema 2.2.7]{Li1}.
\end{example}

% segment-id: o014.li2.chapter7.beck-monadicity-a.p006
Karena endofunktor beraksi pada $\mathcal{C}$, kita ingin membahas ``modul''
di dalam $\mathcal{C}$ di bawah aksi monad $T$. Karena alasan historis,
struktur semacam ini juga disebut aljabar-$T$.

% segment-id: o014.li2.chapter7.beck-monadicity-a.d002
\begin{definition}[S.\ Eilenberg, J.\ C.\ Moore]\label{def:T-Alg}
	\index{modul-T@modul-$T$}
	Misalkan $(T,\mu,\eta)$ monad di atas $\mathcal{C}$. Suatu \emph{modul-$T$}
	adalah data $(M,a)$, dengan $M\in\Obj(\mathcal{C})$ dan $a$ morfisme
	$T(M)\to M$ dalam $\mathcal{C}$, sedemikian sehingga diagram berikut
	komutatif:
	% segment-id: o014.li2.chapter7.beck-monadicity-a.q008
	\[\begin{tikzcd}
		T^2 (M) \arrow[r, "Ta"] \arrow[d, "\mu_M"'] & T(M) \arrow[d, "a"] \\
		T(M) \arrow[r, "a"'] & M
	\end{tikzcd} \quad \begin{tikzcd}
		M \arrow[r, "\eta_M"] \arrow[rd, "{\identity_M}"'] & T(M) \arrow[d, "a"] \\
		& M .
	\end{tikzcd}\]

	Semua modul-$T$ membentuk kategori $\mathcal{C}^T$. Morfisme dari $(M,a)$
	ke $(M',a')$ didefinisikan sebagai morfisme $f:M\to M'$ dalam
	$\mathcal{C}$ yang membuat diagram berikut komutatif:
	% segment-id: o014.li2.chapter7.beck-monadicity-a.q009
	\[\begin{tikzcd}
		T(M) \arrow[r, "a"] \arrow[d, "Tf"'] & M \arrow[d, "f"] \\
		T(M')  \arrow[r, "{a'}"'] & M' .
	\end{tikzcd}\]

	Secara dual, misalkan $(L,\delta,\epsilon)$ komonad di atas $\mathcal{C}$.
	Suatu komodul-$L$ adalah data $(N,b)$, dengan
	$b:N\to L(N)$ suatu morfisme, dan disyaratkan memenuhi diagram komutatif
	yang dual terhadap diagram di atas. Kategori komodul-$L$ ditulis
	$\mathcal{C}^L$.
\end{definition}

% segment-id: o014.li2.chapter7.beck-monadicity-a.p007
Kedua diagram komutatif untuk modul-$T$ masing-masing harus dipandang sebagai
hukum asosiativitas dan hukum unit bagi aksi $T$. Contoh-contoh dan sifat-
sifat berikut terutama dinyatakan untuk monad dan modul di atasnya; versi
komonad dan komodul sepenuhnya dual.

% segment-id: o014.li2.chapter7.beck-monadicity-a.eg002
\begin{example}\label{eg:Mon-monad}
	Tuliskan $\cate{Mon}$ untuk kategori monoid, yang secara konvensi dianggap
	direalisasikan pada himpunan kecil, dan tinjau pasangan adjoin
	% segment-id: o014.li2.chapter7.beck-monadicity-a.q010
	\[\begin{tikzcd}
		\mathbf{M}: \cate{Set} \arrow[r, shift left] & \cate{Mon}: U, \arrow[l, shift left]
	\end{tikzcd}\]
	dengan $\mathbf{M}$ memetakan himpunan $X$ ke monoid bebas
	$\mathbf{M}(X)$, yang definisi dan konstruksinya diberikan dalam
	\cite[Definisi 4.8.1, Lema 4.8.4]{Li1}, dan $U$ adalah funktor pelupa.
	Dengan demikian,
	% segment-id: o014.li2.chapter7.beck-monadicity-a.q011
	\[ T := U\mathbf{M}: \cate{Set} \to \cate{Set}, \quad T(X) = \bigsqcup_{n \geq 0} X^n. \]
	Dengan kata lain, unsur $T(X)$ adalah ``kata''
	$(x_1,\ldots,x_n)$ dengan $n\in\Z_{\geq0}$. Unit pasangan adjoin
	$\eta_X:X\to U\mathbf{M}(X)$ memetakan $x\in X$ ke kata berpanjang $1$,
	$(x)$. Untuk monoid $A$, kounit
	$\varepsilon_A:\mathbf{M}U(A)\to A$ memetakan kata
	$(a_1,\ldots,a_n)$ ke hasil kali $a_1\cdots a_n\in A$. Karena perkalian
	dalam $\mathbf{M}(X)$ adalah konkatenasi kata, pemetaan
	$\mu_X=U\varepsilon\mathbf{M}:T^2(X)\to T(X)$ yang bersesuaian dengan monad
	$T$ (atau, jika dipandang sebagai pemetaan
	$\mathbf{M}(\mathbf{M}(X))\to\mathbf{M}(X)$) meratakan ``kata yang
	terdiri atas kata-kata'', yakni menggabungkan serangkaian daftar, seperti
	% segment-id: o014.li2.chapter7.beck-monadicity-a.q012
	\[ \left( (x_1, \ldots, x_n), (y_1, \ldots, y_m), \ldots \right) \mapsto \left( x_1, \ldots, x_n, y_1, \ldots, y_m, \ldots \right). \]

	Memberikan modul-$T$ $(M,a)$ untuk $T=U\mathbf{M}$ sama artinya dengan
	memberikan $m_1\cdots m_n:=a(m_1,\ldots,m_n)$ untuk setiap kata
	$(m_1,\ldots,m_n)$ pada himpunan $M$, sedemikian sehingga $a(m)=m$ dan
	hukum asosiatif
	% segment-id: o014.li2.chapter7.beck-monadicity-a.q013
	\[ a\left( a(x_1, \ldots, x_n), a(y_1, \ldots, y_m), \ldots \right) = a\left( x_1, \ldots, x_n, y_1, \ldots, y_m, \ldots \right) \]
	berlaku. Singkatnya, modul-$T$ tidak lain adalah monoid, dengan citra kata
	kosong ($n=0$) di bawah $a$ sebagai satuannya. Mudah diperiksa bahwa
	morfisme modul-$T$ juga tidak lain adalah homomorfisme monoid.

	Kasus kategori grup $\cate{Grp}$ sepenuhnya serupa; kuncinya ialah mengganti
	$\mathbf{M}$ dengan funktor grup bebas $\mathbf{F}$. Modul-$T$ yang
	bersesuaian tepat merupakan grup.
\end{example}

% segment-id: o014.li2.chapter7.beck-monadicity-a.eg003
\begin{example}\label{eg:Mod-monad}
	Misalkan $R$ gelanggang dan tinjau pasangan adjoin
	% segment-id: o014.li2.chapter7.beck-monadicity-a.q014
	\[\begin{tikzcd}
		\mathbf{F}: \cate{Set} \arrow[r, shift left] & R\dcate{Mod}: U, \arrow[l, shift left]
	\end{tikzcd}\]
	dengan $\mathbf{F}$ memetakan $X$ ke modul-$R$ kiri bebas
	$R^{\oplus X}$ dan $U$ adalah funktor pelupa. Dari sini diperoleh monad
	$(T,\mu,\eta)$ di atas $\cate{Set}$. Seperti dalam
	Contoh~\ref{eg:Mon-monad}, $T=U\mathbf{F}$ memetakan himpunan $X$ ke
	$R^{\oplus X}$ sebagai suatu himpunan; unsur-unsurnya adalah kombinasi
	linear-$R$ formal hingga dari unsur-unsur $X$,
	$\text{\textquotedblleft} r_1x_1+\cdots+r_nx_n
	\text{\textquotedblright}$. Unit $\eta_X$ memetakan $x\in X$ ke
	$\text{\textquotedblleft}x\text{\textquotedblright}$, sedangkan
	$\mu_X:R^{\oplus(R^{\oplus}X)}\to R^{\oplus X}$ meratakan kombinasi linear
	dua tingkat, yakni
	% segment-id: o014.li2.chapter7.beck-monadicity-a.q015
	\begin{multline*}
		\text{\textquotedblleft} \alpha (\text{\textquotedblleft} r_1 x_1 + \cdots + r_n x_n \text{\textquotedblright}) + \beta (\text{\textquotedblleft} s_1 y_1 + \cdots + s_m y_m \text{\textquotedblright}) + \cdots \text{\textquotedblright} \\
		\mapsto \text{\textquotedblleft} (\alpha r_1) x_1 + \cdots + (\alpha r_n) x_n + (\beta s_1) y_1 + \cdots + (\beta s_m) y_m + \cdots \text{\textquotedblright} .
	\end{multline*}
	Memberikan modul-$T$ $(M,a)$ sama artinya dengan memberikan himpunan $M$
	beserta aturan yang memetakan kombinasi linear formal ke suatu unsur $M$,
	memetakan $\text{\textquotedblleft}x\text{\textquotedblright}$ ke $x$,
	bersifat asosiatif terhadap penjumlahan dan perkalian skalar, dan membuat
	perkalian skalar oleh $1\in R$ menjadi pemetaan identitas. Singkatnya,
	modul-$T$ tidak lain adalah modul-$R$ kiri. Jelas bahwa morfisme modul-$T$
	tepat merupakan homomorfisme modul.
\end{example}

% segment-id: o014.li2.chapter7.beck-monadicity-a.p008
Kembali ke teori umum. Terdapat funktor pelupa
$U^T:\mathcal{C}^T\to\mathcal{C}$ yang memetakan objek $(M,a)$ ke $M$.
Berikut diberikan adjoin kirinya: konstruksi bebas.

% segment-id: o014.li2.chapter7.beck-monadicity-a.dp001
\begin{definition-proposition}[Modul-$T$ bebas]\label{def:free-T-Alg}
	\index{modul!bebas@bebas (free)}
	Misalkan $(T,\mu,\eta)$ monad di atas $\mathcal{C}$. Definisikan funktor
	% segment-id: o014.li2.chapter7.beck-monadicity-a.q016
	\[ \mathrm{Free}^T: \mathcal{C} \to \mathcal{C}^T , \quad
	\begin{array}{rl}
		(\text{objek } M) & \mapsto (TM, \mu_M: T^2 M \to TM), \\
		(\text{morfisme } f: M \to M') & \mapsto Tf: TM \to TM'.
	\end{array}\]
	Hal ini memberikan pasangan adjoin
	% segment-id: o014.li2.chapter7.beck-monadicity-a.q017
	\[\begin{tikzcd}
		\mathrm{Free}^T : \mathcal{C} \arrow[shift left, r] & \mathcal{C}^T: U^T \arrow[shift left, l],
	\end{tikzcd}\]
	dan monad pada $\mathcal{C}$ yang ditentukan pasangan adjoin ini tepat
	$(T,\mu,\eta)$.
\end{definition-proposition}
% segment-id: o014.li2.chapter7.beck-monadicity-a.pf001
\begin{proof}
	Untuk memeriksa bahwa $(TM,\mu_M)$ adalah modul-$T$, perlu ditunjukkan
	komutativitas kedua diagram berikut:
	% segment-id: o014.li2.chapter7.beck-monadicity-a.q018
	\[\begin{tikzcd}
		T^3 (M) \arrow[r, "T\mu_M"] \arrow[d, "\mu_{T(M)}"'] & T^2(M) \arrow[d, "\mu_M"] \\
		T^2 (M) \arrow[r, "\mu_M"'] & T(M)
	\end{tikzcd} \quad \begin{tikzcd}
		T(M) \arrow[r, "\eta_{T(M)}"] \arrow[rd, "\identity"'] & T^2(M) \arrow[d, "\mu_M"] \\
		& T(M).
	\end{tikzcd}\]
	Namun, diagram ini diperoleh dengan mengevaluasi diagram funktor dalam
	Definisi~\ref{def:monad} pada objek $M$. Funktorialitasnya jelas.

	Selanjutnya, perhatikan bahwa $U^T\mathrm{Free}^T=T$. Ambil
	$\eta:\identity_{\mathcal{C}}\to T$ yang merupakan bagian dari data $T$, lalu
	definisikan
	% segment-id: o014.li2.chapter7.beck-monadicity-a.q019
	\[ \epsilon: \mathrm{Free}^T U^T \to \identity_{\mathcal{C}^T}, \quad
	\epsilon_{(M, a)}: (TM, \mu_M) \xrightarrow{a} (M, a) \in \Obj(\mathcal{C}^T). \]
	Diagram kiri dalam Definisi~\ref{def:T-Alg} menunjukkan bahwa
	$\epsilon_{(M,a)}$ adalah morfisme dalam $\mathcal{C}^T$.
	Funktorialitasnya juga jelas.

	Pemeriksaan rutin menunjukkan bahwa $(\mathrm{Free}^T,U^T)$ membentuk
	pasangan adjoin dengan unit $\eta$ dan kounit $\epsilon$; perinciannya
	dihilangkan. Selain itu,
	% segment-id: o014.li2.chapter7.beck-monadicity-a.q020
	\begin{equation*}\begin{aligned}
		\left(U^T \epsilon \mathrm{Free}^T\right)_M & = U^T \epsilon_{(TM, \mu_M)} \\
		& = \left[ \mu_M: T^2(M) \to T(M)\right],
	\end{aligned}\end{equation*}
	yang menunjukkan bahwa monad di atas $\mathcal{C}$ yang ditentukan oleh
	pasangan adjoin $(\mathrm{Free}^T,U^T)$ adalah $(T,\mu,\eta)$.
\end{proof}

% segment-id: o014.li2.chapter7.beck-monadicity-a.p009
Untuk monad dalam Contoh~\ref{eg:Mon-monad} dan
\ref{eg:Mod-monad}, funktor $\mathrm{Free}^T$ masing-masing bersesuaian dengan
konstruksi monoid bebas dan modul bebas.

% segment-id: o014.li2.chapter7.beck-monadicity-a.p010
Secara dual, komonad $L$ di atas kategori $\mathcal{D}$ juga memberikan
pasangan adjoin bebas--pelupa; tidak perlu dibahas lebih lanjut.
